\subsection{The exceptional center at the remaining binary-arm stars}

\begin{proposition}[Proved: a binary arm forces the two-by-three band center]
\label{CP2:A-star-binary-arm-band-center}
Let $C$ be a source-star category with at most eighteen
arrows.  Suppose its center monoid $M$ has order at most
seven, each leaf monoid has order at most four, and at
least one stalk has at most two arrows.  If
$BC\simeq\CP^2$, then
\[
 M\cong (\{0,1\}\times\{0,1,2\})^1.
\]
Moreover, each binary stalk has two states and its right
$M$ action is a surjective quotient of the three labelled
minimal columns onto those two states.  The corresponding
leaf monoid acts trivially on that stalk.
\end{proposition}
\begin{proof}
Both leaf corners have zero rational second cohomology
by the small-minimal-ideal theorem.  Write $S$ for a stalk
with at most two states, $N$ for its leaf monoid, and
$E_S=B(N,S,M)$ for its bridge.

If $H^2(BM;\mathbb Q)=0$, all corner restrictions of a
global second class vanish.  The two-block relative-cup
lemma, with the source corner as one block and the
disjoint union of the leaves as the other, makes every
second square zero.  We may therefore assume
$H^2(BM;\mathbb Q)\ne0$.

If the right $M$ act $S$ has a fixed state, its inclusion
gives a section $BM\to B(S,M)$ of the action projection.
The inclusion $B(S,M)\to E_S$ uses empty $N$ strings.
On the bridge a global second class has equal pullbacks
from the two corners, while its restriction to $BN$ is
zero.  Restriction along this section consequently makes
its restriction to $BM$ zero.  Again the two-block
relative-cup lemma kills its square.

If every right actor on $S$ is a permutation, the
projection $B(S,M)\to BM$ is a nonempty finite covering.
Rational transfer makes its cohomology pullback injective.
The same compatibility of bridge restrictions therefore
makes every global second class restrict to zero on
$BM$, and the preceding relative-cup argument applies.
These cases include every singleton stalk.

The remaining action has exactly two states, no fixed
state, and a nonpermutation actor.  A nonpermutation of
two points is constant.  If only one of the two constant
maps occurred, every actor would fix its value: a
transposition would produce the other constant by
composition.  Thus both constant maps occur.  The image
of the source's minimal ideal is the two-constant ideal
of this transformation monoid.

The nonzero second cohomology forces at least two
minimal rows and columns, by
Corollary~\ref{A:monoids-nine}.  The order-seven budget
forces the minimal group to be trivial, so the minimal
ideal $K$ is a rectangular band with row and column
counts $r,c\ge2$.  Its action on $S$ sends each labelled
column to one of the two states, independently of the
row, and this column map is onto.  One direct check is
to write $\rho_{ij}$ for a right band actor.  Since the
image is the two-constant ideal, each $\rho_{ij}$ is a
constant, and
$\rho_{kl}\rho_{ij}=\rho_{il}$ makes its value depend
only on $j$.

If $c=2$, the column map is a bijection.  For any
minimal idempotent $e$, the right ideal $eM$ consists
of exactly those two column elements, and the preceding
identification is a literal right-$M$-act isomorphism
$S\cong eM$.  Its action category is contractible:
$eM$ is a retract of the regular right act $M$, whose
action category has an initial object.  A transformation
of $S$ commuting with both constants is the identity,
so every left $N$ actor is trivial.  Therefore
\[
 E_S=BN\times B(eM,M)\longrightarrow BN
\]
is a homotopy equivalence.  In the homotopy pushout
$BM\leftarrow E_S\to BN$, the inclusion of $BM$ into
this entire edge category is a homotopy equivalence.
The full star is obtained by gluing its other edge along
that same $BM$.  Since the nerve inclusions are CW
cofibrations, deleting the binary edge and its leaf
preserves the classifying-space homotopy type.  The
remaining two-object category has at most
$18-(|N|+2)\le15$ arrows, contradicting
Theorem~\ref{CP2:A-global-lower-eighteen}.

We must have $c\ge3$.  Now $r,c\ge2$ and
$rc+1\le7$ force $r=2$, $c=3$, and only the global
identity outside $K$.  This is exactly the stated
seven-element monoid.  The already proved description
of the binary action and its commuting left action
gives the remaining assertions.
\end{proof}

\begin{corollary}[Proved: the exact center scope for three pending star profiles]
\label{CP2:A-star-three-binary-center-scope}
For the original star profiles
\[
 52=(7,1,3,2,5),\qquad
 53=(7,1,4,2,4),\qquad
 60=(7,3,3,2,3),
\]
only the single center isomorphism class
$(\{0,1\}\times\{0,1,2\})^1$ can occur in a
projective-plane model.  This is an analytic restriction
of the center-table scope; it does not complete the
remaining action and composition checks.
\end{corollary}
\begin{proof}
Every displayed profile satisfies all hypotheses of
the proposition, using its first stalk as the binary arm.
\end{proof}
